% Part: sets-functions-relations
% Chapter: relations
% Section: special-properties
% Hindi translation (hi-Deva-IN) of Open Logic commit 9620cc73f9c8e0ad003c514a5d3748f29611c4c0.

\documentclass[../../../include/open-logic-section]{subfiles}

\begin{document}

\olfileid[hi]{sfr}{rel}{prp}
\olsection{संबंधों के विशेष गुणधर्म}

\begin{intro}
कुछ प्रकार के संबंध इतने सामान्य पाए जाते हैं कि उन्हें विशेष नाम दिए गए हैं।
उदाहरण के लिए, $\le$ और~$\subseteq$ दोनों अपने-अपने प्रांतों के अवयवों को
समान ढंग से संबंधित करते हैं---इनके प्रांत क्रमशः $\Nat$ (संबंध~$\le$ के
लिए) और $\Pow{A}$ (संबंध~$\subseteq$ के लिए) हैं। यह ठीक-ठीक समझने के लिए कि ये संबंध किस प्रकार
समान हैं और किस प्रकार भिन्न, हम उन विशेष गुणधर्मों के अनुसार उनका वर्गीकरण
करते हैं जो संबंधों में हो सकते हैं। इन विशेष गुणधर्मों में से कुछ (और उनके
संयोजन) विशेष रूप से महत्त्वपूर्ण हैं: क्रम और तुल्यता संबंध।
\end{intro}

\begin{defn}[स्वतुल्यता]
कोई संबंध $R \subseteq A^2$ \emph{स्वतुल्य} होता है तभी और केवल तभी जब
प्रत्येक $x \in
A$ के लिए $Rxx$ हो।
\end{defn}

\begin{defn}[संक्रामकता]
कोई संबंध $R \subseteq A^2$ \emph{संक्रामक} होता है तभी और केवल तभी जब
$Rxy$ और $Ryz$ होने पर $Rxz$ भी हो।
\end{defn}

\begin{defn}[सममिति]
कोई संबंध~$R \subseteq A^2$ \emph{सममित} होता है तभी और केवल तभी जब
$Rxy$ होने पर~$Ryx$ भी हो।
\end{defn}

\begin{defn}[प्रति-सममिति]
कोई संबंध~$R \subseteq A^2$ \emph{प्रति-सममित} होता है तभी और केवल तभी जब
$Rxy$ और $Ryx$ दोनों का सत्य होना $x=y$ को अनिवार्य करे (या दूसरे शब्दों में: यदि $x\neq y$
हो, तो या तो $\lnot Rxy$ अथवा $\lnot Ryx$ हो)।
\end{defn}

\begin{explain}
सममित संबंध में $Rxy$ और $Ryx$ या तो सदैव साथ-साथ सत्य होते हैं, अथवा
दोनों में से कोई भी सत्य नहीं होता। प्रति-सममित संबंध में $Rxy$ और $Ryx$
दोनों साथ-साथ केवल तभी सत्य हो सकते हैं जब $x = y$ हो। ध्यान दें कि इससे यह
\emph{आवश्यक} नहीं होता कि $Rxy$ और $Ryx$, $x = y$ होने पर सत्य हों; केवल इतना है
कि इसकी मनाही नहीं है। अतः कोई प्रति-सममित संबंध स्वतुल्य हो सकता है, परंतु
हर प्रति-सममित संबंध स्वतुल्य नहीं होता। यह भी ध्यान दें कि प्रति-सममित
होना और केवल सममित न होना अलग-अलग शर्तें हैं। वास्तव में कोई संबंध एक ही
समय में सममित और प्रति-सममित दोनों हो सकता है (जैसे तत्समक संबंध)।
\end{explain}

\begin{defn}[संबद्धता]
कोई संबंध $R \subseteq A^2$ \emph{संबद्ध} होता है यदि सभी $x,y\in
A$ के लिए, $x \neq y$ होने पर या तो $Rxy$ अथवा~$Ryx$ हो।
\end{defn}

\begin{prob}
ऐसे संबंधों के उदाहरण दीजिए जो (a) स्वतुल्य और सममित हों, परंतु संक्रामक
न हों; (b) स्वतुल्य और प्रति-सममित हों; (c) प्रति-सममित और संक्रामक हों,
परंतु स्वतुल्य न हों; तथा (d) स्वतुल्य, सममित और संक्रामक हों। संख्याओं या
समुच्चयों पर बने संबंधों का उपयोग न कीजिए।
\end{prob} 
 
\begin{defn}[अस्वतुल्यता]
कोई संबंध $R \subseteq A^2$ \emph{अस्वतुल्य} कहलाता है यदि प्रत्येक $x \in
A$ के लिए $Rxx$ न हो। 
\end{defn}

\begin{defn}[असममिति]
कोई संबंध $R \subseteq A^2$ \emph{असममित} कहलाता है यदि ऐसा कोई युग्म $x,y\in
A$ न हो जिसके लिए $Rxy$ और~$Ryx$ दोनों हों। 
\end{defn}

ध्यान दें कि यदि $A \neq \emptyset$ हो, तो~$A$ पर कोई भी अस्वतुल्य संबंध
स्वतुल्य नहीं होता, और~$A$ पर प्रत्येक असममित संबंध प्रति-सममित भी होता है।
परंतु ऐसे $R \subseteq A^2$ होते हैं जो न स्वतुल्य होते हैं और न अस्वतुल्य,
और ऐसे प्रति-सममित संबंध भी होते हैं जो असममित नहीं होते। 

\end{document}
